% =====================================================================
\section{Risks, Limitations and Mitigation}
\label{sec:risks}
% =====================================================================

\Cref{tab:risks} lists the main implementation risks, an assessment of their
likelihood and impact, and the planned mitigation. Assessments are the team's
judgement at the planning stage and will be reviewed weekly.

{\small
\setlength{\tabcolsep}{4pt}
\renewcommand{\arraystretch}{1.15}
\begin{longtable}{@{}p{3.5cm}p{1.4cm}p{10.2cm}@{}}
\caption{Implementation risks and mitigation strategies (L/I = likelihood/impact, rated H = high, M = medium, L = low).}\label{tab:risks}\\
\toprule
\textbf{Risk} & \textbf{L / I} & \textbf{Mitigation} \\
\midrule
\endfirsthead
\multicolumn{3}{@{}l}{\small\itshape \Cref{tab:risks} continued}\\
\toprule
\textbf{Risk} & \textbf{L / I} & \textbf{Mitigation} \\
\midrule
\endhead
\bottomrule
\endfoot
Limited compute resources & H / H &
Confirm GPU access in week~1; use mixed precision and cached crops; develop on a
subset of speakers; keep the model moderate in size; drop optional experiments
first (\cref{sec:schedule}). \\
Long training times & H / M &
Start B1 in week~2 and overlap training with development; use early stopping;
schedule long runs overnight; checkpoint every epoch so runs can be resumed. \\
Misaligned audio and video & M / H &
Never drop frames during cropping; integrity check of durations; automated
train/infer parity test; time-shift sanity test (\cref{tab:tests}); clap-test
calibration for live input. \\
CTC training instability & M / H &
Learning-rate warm-up, gradient clipping and \texttt{zero\_infinity}; assert the
length condition (\cref{eq:ctc-cond}); overfit a tiny batch before full runs;
log non-finite losses. \\
Token distribution issues & M / L &
The blank symbol dominates CTC outputs and GRID's grammar makes some characters
much more frequent; monitor CER per word slot and the share of blank predictions;
inspect the letter slot separately, as it is the hardest. \\
Speaker leakage & L / H &
Speaker-level split stored in one file; loader refuses overlapping speakers;
disjointness unit test; statistics from training speakers only. \\
Poor generalisation to unseen speakers & M / H &
Crop jitter, horizontal flip, dropout and weight decay; model selection on
held-out validation speakers; report per-speaker results honestly. \\
Mouth detection failure (lighting, head pose, occlusion) & M / M &
GRID is frontal and well lit, so this mainly affects the live demo; interpolate
short gaps; fall back to audio-only captions using $\mathbf{e}_{\text{miss}}$; good
lighting and frontal camera position for the demonstration. \\
Attention fusion fails to improve performance & M / H &
Report the outcome truthfully; diagnose with gate analysis, test-time modality
ablations and the checklist in \cref{sec:evaluation}; tune modality dropout;
B3 remains a working fused system. \\
Excessive inference latency & M / M &
Measure early on the demo machine; increase window hop, reduce beam width, use a
GPU or quantise; Mode A batch captioning as the guaranteed fallback. \\
Limited multilingual generalisation & H / L (core) &
Treat language identification as an optional extension with its own data; state
clearly that the core model covers GRID English only. \\
Domain gap between GRID and live speech & H / M &
Restrict the demonstration to GRID-grammar sentences; explain the closed
vocabulary; identify larger datasets (e.g.\ those in~\cite{afouras2022}) as future
work. \\
Privacy of camera and microphone input & M / H &
Process all input locally; do not store live video or audio by default; show a
clear on-screen indicator while capturing; record any demonstration data only
with written consent and delete it after the project; never upload live input to
third-party services. \\
\end{longtable}
}

\subsection{Limitations of the planned system}
Even if every risk is managed, the system will have known limitations. GRID's
fixed grammar and 51-word vocabulary~\cite{cooke2006} make the recognition task far
easier than open captioning, so a low WER on GRID does not imply comparable
accuracy on lectures or meetings~\cite{afouras2022}. Synthetic additive noise does
not reproduce reverberation, competing speakers or microphone distortion. The
visual front-end assumes a roughly frontal, unoccluded face, which the literature
identifies as an open problem~\cite{ameer2025,kanavnet2026}. The final report will
state these limitations alongside the results.
